% =====================================================================
% AVISO DE VIGENCIA -- 15-08-2026
%
% ESTE ANEXO NO SE COMPILA EN `main.tex` Y ESTA SUPERADO EN PARTE.
%
% Describe el flujo estadistico tal como era ANTES de las remediaciones
% P0-B, P0-C, P0-G y P0-H. En concreto, ya no son vigentes:
%
%   * la activacion de modelos "por niveles" (retirada en P0-B);
%   * los umbrales de ventanas 6 / 3 usados como VETO: desde P0-G no vetan,
%     se comunican como advertencia y no cancelan el pronostico puntual;
%   * la regla de prefijo -- "un horizonte solo es defendible si todos los
%     anteriores lo son" -- retirada en P0-H: ninguna fuente exige que los
%     horizontes validos formen un prefijo;
%   * cualquier mencion a "proyeccion tecnica validada" o a "certificar" un
%     horizonte: no existe ninguna validacion en ese sentido;
%   * la publicacion del intervalo de prediccion, retirada en P0-C / C2, junto
%     con su ancho relativo, su cobertura y los tres tramos que aparecen mas
%     abajo: ninguno clasifica ni bloquea nada;
%   * el AJUSTE de cambio de anio (gamma y el factor f_h), retirado en P0-F el
%     12-08-2026: `aplicado = False` de forma incondicional. El patron se sigue
%     midiendo y publicando como diagnostico descriptivo, pero NO modifica el
%     pronostico. Su incorporacion predictiva queda como trabajo futuro.
%
% El estado vigente esta en el capitulo correspondiente de `main.tex` y en
% `tabla_criterios_auditable.tex`. Este anexo se conserva como registro
% historico del diseno anterior; no debe citarse como descripcion actual.
% =====================================================================

% Flujo estadistico de la aplicacion ICOCIV, paso a paso.
% Documento autocontenido: compila por si solo con pdflatex.
% Para integrarlo en el documento final, copie o \input el contenido que va
% desde \section{Proposito} hasta el final, omitiendo el preambulo.
\documentclass[11pt,letterpaper]{article}
\usepackage[utf8]{inputenc}
\usepackage[T1]{fontenc}
\usepackage[spanish,es-noquoting]{babel}
\usepackage[margin=2.5cm]{geometry}
\usepackage{amsmath}
\usepackage{booktabs}
\usepackage{longtable}
\usepackage{array}
\usepackage{hyperref}

\title{Flujo estadístico de la aplicación ICOCIV\\\large Explicación paso a paso}
\author{Proyecto ICOCIV / ICCP--ICOCIV}
\date{\today}

\begin{document}
\maketitle

\section{Propósito}

Este anexo explica, paso a paso, qué hace la aplicación en su componente
estadístico: desde la selección de una ruta jerárquica ICOCIV hasta la
generación, advertencia o negación de una proyección. El objetivo es que un
lector técnico pueda auditar cada decisión y reproducirla.

La aplicación \textbf{apoya} el criterio técnico del profesional responsable;
\textbf{no lo reemplaza}. Ninguna salida sustituye la validación del ingeniero
ni la producción estadística oficial del DANE.

\section{Pasos del flujo}

\begin{longtable}{@{}p{0.5cm}p{4.6cm}p{9.4cm}@{}}
\toprule
\textbf{\#} & \textbf{Paso} & \textbf{Qué hace la aplicación} \\
\midrule
\endhead
1 & Selección de ruta ICOCIV & El usuario recorre grupo de obra, subclase CPC, tipología, capítulo, grupo de costos e insumo. La ruta determina la tabla y la fila fuente del anexo. \\
2 & Extracción de la serie & Se localiza la fila fuente y se leen las columnas de periodos mensuales, conservando la trazabilidad de origen. \\
3 & Normalización temporal & Los periodos se convierten a una secuencia mensual ordenada (AAAA--MM) y se ordena cronológicamente. \\
4 & Validación de calidad & Se revisan longitud, continuidad, duplicados, faltantes, valores no numéricos y valores no positivos. \\
5 & Análisis descriptivo & Media, mediana, mínimo, máximo, desviación y coeficiente de variación del nivel del índice. \\
6 & Variaciones & Variación mensual y variación acumulada entre extremos de la serie. \\
7 & Detección de atípicos & Puntaje modificado con mediana y MAD sobre las variaciones. Se \emph{reportan}; no se eliminan ni se imputan valores oficiales. \\
8 & Modelos candidatos & Activación progresiva por niveles según longitud, horizonte y diagnóstico. \\
9 & Backtesting walk-forward & Ventana expansiva: en cada origen se reentrena solo con datos anteriores y se pronostica el horizonte objetivo. \\
10 & Métricas de error & MAE, RMSE, MAPE, sMAPE, MASE, sesgo medio, estabilidad y porcentaje de errores extremos, por modelo y por horizonte. \\
11 & Comparación de modelos & Cada candidato se contrasta contra los benchmarks Naive y Drift mediante errores relativos. \\
12 & Horizonte recomendado & Último horizonte consecutivo clasificado como proyección técnica (terminología anterior a P0-G; no implica validación). \\
13 & Horizonte admisible & Último horizonte consecutivo aceptable al menos como escenario. \\
14 & Decisión & Se genera la proyección, se genera como escenario con advertencia, o se niega indicando la causa. \\
15 & Presentación & Tarjetas de resultado, tabla por horizonte, serie, fila fuente y gráfica interactiva. \\
16 & Exportación & Informe DOCX y PDF, CSV reproducible y Excel acumulativo, con la misma trazabilidad. \\
\bottomrule
\end{longtable}

\section{Fórmulas principales}

\subsection{Variaciones}
\begin{equation}
v_t=\frac{y_t-y_{t-1}}{y_{t-1}}\times 100,
\qquad
v_{\text{acum}}=\frac{y_T-y_1}{y_1}\times 100 .
\end{equation}

\subsection{Detección robusta de atípicos}
Con mediana \(\tilde{v}\) y desviación absoluta mediana \(\mathrm{MAD}\):
\begin{equation}
M_i=\frac{0{,}6745\,(v_i-\tilde{v})}{\mathrm{MAD}},
\qquad \text{se marca si } |M_i|>3{,}5 .
\end{equation}

\subsection{Backtesting walk-forward}
Para cada origen \(r\) y horizonte \(h\), el modelo se reestima con
\(\{y_1,\dots,y_r\}\) y se compara \(\hat y_{r+h|r}\) contra \(y_{r+h}\):
\begin{equation}
e_{r,h}=y_{r+h}-\hat y_{r+h|r}.
\end{equation}
Este orden garantiza que nunca se use información futura.

\subsection{Métricas por horizonte}
Con \(m\) orígenes válidos:
\begin{align}
\mathrm{MAE}_h&=\frac{1}{m}\sum_r |e_{r,h}|, &
\mathrm{RMSE}_h&=\sqrt{\frac{1}{m}\sum_r e_{r,h}^{2}},\\
\mathrm{MAPE}_h&=\frac{100}{m}\sum_r\left|\frac{e_{r,h}}{y_{r+h}}\right|, &
\mathrm{sMAPE}_h&=\frac{100}{m}\sum_r\frac{2|e_{r,h}|}{|y_{r+h}|+|\hat y_{r+h|r}|}.
\end{align}
El MASE escala cada error con la referencia \emph{naive} calculada
\textbf{dentro de la ventana de entrenamiento} de cada origen, de modo que la
escala nunca incorpora información del tramo evaluado:
\begin{equation}
\mathrm{MASE}_h=\frac{1}{m}\sum_r \frac{|e_{r,h}|}{s_r},
\qquad
s_r=\frac{1}{r-1}\sum_{t=2}^{r}|y_t-y_{t-1}| .
\end{equation}

\subsection{Selección del modelo de la trayectoria}
\label{sec:modelo}

El backtesting produce un modelo ganador \emph{para cada horizonte}. Si la
trayectoria se generara con el ganador del horizonte solicitado, los primeros
meses cambiarían según lo que pidiera el usuario: para una misma serie, el mes
siguiente al último dato podía diferir en más de un 2\,\% entre una solicitud a
3 meses y otra a 18. Eso es indeseable, porque el ajuste al comportamiento
histórico no depende de hasta dónde se quiera proyectar.

Por eso se elige \textbf{un único modelo por serie} para toda la trayectoria.
Cada modelo \(k\) se puntúa por su RMSE relativo al mejor de cada horizonte,
ponderando cada horizonte por \(1/h\):

\begin{equation}
S_k=\frac{\displaystyle\sum_{h}\frac{1}{h}\cdot\frac{\mathrm{RMSE}_{k,h}}{\min_j \mathrm{RMSE}_{j,h}}}
{\displaystyle\sum_{h}\frac{1}{h}},
\qquad
k^{*}=\arg\min_k S_k .
\end{equation}

La ponderación \(1/h\) refleja que el error crece con el horizonte y que la
evidencia de corto plazo es más confiable y más consultada; sin ella, un modelo
mediocre pero parejo en horizontes largos puede desplazar a otro que domina en
el tramo corto. Solo compiten modelos evaluados en todos los horizontes.

El conjunto de horizontes evaluados depende únicamente de la longitud de la
serie, de modo que \(k^{*}\) ---y por tanto los primeros valores
proyectados--- es una propiedad de la serie y no de la solicitud del usuario.

El modelo \(k^{*}\) se fija \emph{antes} de evaluar los horizontes, de manera
que la clasificación de admisibilidad, las métricas reportadas, los intervalos
y la trayectoria correspondan todos al mismo modelo. La admisibilidad se sigue
decidiendo horizonte por horizonte, pero siempre con \(k^{*}\): así se evita
informar un modelo en la trayectoria y sustentar la advertencia en otro.

La combinación ponderada de varios modelos se evaluó y se descartó: sus pesos
por horizonte reintroducían la dependencia respecto del horizonte solicitado, y
la variante de pesos fijos por serie no mejoró al modelo único al medirla sin
fuga de información.

\subsection{Comparación contra benchmarks}
\begin{equation}
\mathrm{rRMSE}_{\text{naive}}=\frac{\mathrm{RMSE}_{\text{modelo}}}{\mathrm{RMSE}_{\text{naive}}},
\end{equation}
y de forma análoga frente a Drift. Un valor menor o igual que 1 indica que el
modelo iguala o mejora la referencia simple.

\subsection{Intervalos de predicción}
Se priorizan los errores fuera de muestra del mismo horizonte, usando
percentiles empíricos (2{,}5/97{,}5 para el intervalo del 95\%). Cuando la
evidencia empírica no alcanza, se emplean cuantiles normales
(\(1{,}2816\) y \(1{,}96\)). El ancho relativo del IC95 respecto al valor
proyectado participa en la clasificación del horizonte.

\section{Modelos candidatos}

La aplicación evalúa únicamente modelos implementados y verificables:

\begin{itemize}
  \item \textbf{Benchmarks:} Naive, Drift, promedio móvil y variación reciente.
  \item \textbf{Suavizamiento:} Holt lineal y Holt con tendencia amortiguada, ambos con coeficientes fijos declarados. ETS no forma parte del catálogo de modelos ejecutados; la decisión de alcance se documenta en el capítulo de pruebas.
  \item \textbf{Regresiones interpretables:} lineal, logarítmica temporal, exponencial/log-lineal y robusta de Huber.
  \item \textbf{Modelos sobre variación:} variación mensual y log-variación.
\end{itemize}

La combinación ponderada de modelos fue evaluada y \textbf{descartada}; la
proyección usa siempre un único modelo por serie (ver \S\ref{sec:modelo}).

La activación es progresiva: el nivel 1 (benchmarks y Holt) siempre está
disponible y el nivel 2 se habilita según longitud, horizonte, volatilidad y
presencia de atípicos. Un modelo que falla al ajustarse se descarta con su
motivo registrado y no contamina la decisión final.

\section{Criterios de decisión}

\begin{longtable}{@{}p{4.3cm}p{2.4cm}p{7.8cm}@{}}
\toprule
\textbf{Criterio} & \textbf{Umbral} & \textbf{Efecto} \\
\midrule
\endhead
Longitud mínima para modelar & 8 observaciones & Por debajo no se modela. \\
Advertencia por longitud & 18 observaciones & Se modela con advertencia. \\
Longitud recomendada & 24 observaciones & Referencia de buena práctica. \\
Primer origen de la validación temporal & \(N_0=\max_m N_{\min}(m)\) (hoy 6) & Define cuántos orígenes quedan disponibles. Corresponde al máximo de los mínimos operativos codificados; no es un mínimo de identificabilidad demostrado. Se deriva de la estimabilidad de los modelos, no de una proporción. \\
Ventanas para \textbf{proyección técnica} & \(\geq 6\) & Umbral anterior a P0-G; desde entonces no veta, se comunica como advertencia. \\
Ventanas para \textbf{escenario} & \(\geq 3\) y \(< 6\) & Evidencia reducida: se calcula con advertencia. Desde P0-G no impide publicar el pronóstico puntual. \\
Ventanas insuficientes & \(< 3\) & No se evalúa el horizonte; se niega indicando la causa. \\
Atípico robusto & \(|M_i|>3{,}5\) & Se reporta como advertencia; no bloquea por sí solo. \\
MAPE extremo & \(> 25\%\) & Bloquea el horizonte. \\
sMAPE extremo & \(> 30\%\) & Bloquea el horizonte. \\
Errores extremos recurrentes & \(\geq 50\%\) de ventanas & Bloquea el horizonte. \\
Desempeño frente a benchmark & \(\mathrm{rRMSE}>1{,}25\) & Bloquea el horizonte. \\
Ancho relativo IC95 & Ver cuadro & Tres tramos por horizonte: aceptable, escenario con advertencia y bloqueo. \\
MASE \(>1\) & --- & Advertencia auxiliar; no bloquea por sí sola. \\
\bottomrule
\end{longtable}

\subsection{Amplitud de los intervalos por horizonte}

El ancho relativo del IC95 respecto al valor proyectado se compara con tres
umbrales que crecen con el horizonte, porque la incertidumbre legítima aumenta
al alejarse en el tiempo:

\begin{center}
\begin{tabular}{@{}cccc@{}}
\toprule
\textbf{Horizonte} & \textbf{Aceptable} & \textbf{Escenario} & \textbf{Bloqueo} \\
\midrule
\(h\leq 3\)  & \(\leq 10\%\) & \(10\%\)--\(30\%\) & \(>30\%\) \\
\(h\leq 6\)  & \(\leq 20\%\) & \(20\%\)--\(40\%\) & \(>40\%\) \\
\(h\leq 12\) & \(\leq 30\%\) & \(30\%\)--\(50\%\) & \(>50\%\) \\
\(h\leq 18\) & \(\leq 40\%\) & \(40\%\)--\(55\%\) & \(>55\%\) \\
\(h> 18\)    & \(\leq 45\%\) & \(45\%\)--\(75\%\) & \(>75\%\) \\
\bottomrule
\end{tabular}
\end{center}

El tramo intermedio es esencial: un horizonte con intervalos amplios pero no
desproporcionados se ofrece como \emph{escenario} con advertencia, en lugar de
bloquearse. Sin ese tramo, superar el umbral por un margen pequeño anulaba el
horizonte y, por el encadenamiento descrito abajo, también todos los
posteriores.

\subsection{Horizonte recomendado y horizonte admisible}

Ambos se calculan recorriendo los horizontes de forma \emph{consecutiva} desde
\(h=1\), porque un horizonte solo es defendible si todos los anteriores lo son:

\begin{itemize}
  \item \textbf{Máximo recomendado:} último \(h\) consecutivo clasificado como proyección técnica (terminología anterior a P0-G).
  \item \textbf{Máximo admisible:} último \(h\) consecutivo aceptable al menos como escenario.
\end{itemize}

No son sinónimos y la aplicación los reporta por separado. Los accesos rápidos
de la interfaz (1, 3, 6, 12, 18 meses) son atajos de solicitud y no constituyen
un límite estadístico.

\section{Ejemplo conceptual con horizonte de 3 meses}

Sea una serie mensual con \(n=24\) observaciones. El primer origen no se fija
con una proporción sino que se \emph{deriva}: es el mayor de los mínimos de
observaciones que exigen los modelos que compiten. Con el catálogo actual ese
máximo lo impone Holt con tendencia amortiguada, cuyos cinco parámetros
---\(\alpha\), \(\beta^{*}\), \(\varphi\), \(\ell_0\) y \(b_0\)--- quedan
identificados a partir de la sexta observación, de modo que \(N_0=6\). El
número de orígenes disponibles para un horizonte \(h\) es entonces
\begin{equation}
m(h)=n-N_0-h+1=24-6-h+1=19-h .
\end{equation}

\begin{center}
\begin{tabular}{@{}ccl@{}}
\toprule
\textbf{h} & \textbf{Orígenes \(m(h)\)} & \textbf{Clasificación} \\
\midrule
1 & 18 & Proyección técnica (\(m\geq 6\), criterio anterior a P0-G) \\
3 & 16 & Proyección técnica (\(m\geq 6\), criterio anterior a P0-G) \\
13 & 6 & Proyección técnica (\(m\geq 6\), criterio anterior a P0-G) \\
15 & 4 & Escenario con advertencia (\(3\leq m<6\)) \\
17 & 2 & No evaluable: sin ventanas suficientes \\
\bottomrule
\end{tabular}
\end{center}

Con este esquema, una solicitud de 3 meses \textbf{sí produce} una proyección.
Una solicitud de 17 meses se niega sobre una serie de 24 observaciones, y el
mensaje indica la causa real: no existen ventanas de validación suficientes, no
un mal desempeño del modelo.

Conviene señalar el efecto del cambio de criterio. Con la regla anterior
---una proporción del 60\,\% con mínimo de 18 observaciones, ninguno de los dos
valores respaldado por fuente--- esa misma serie de 24 observaciones dejaba
\(m(h)=7-h\) orígenes y negaba ya el horizonte de 5 meses. Derivar el primer
origen de la estimabilidad no relaja ningún umbral de error: \textbf{aumenta la
evidencia disponible} y, con ella, el número de horizontes que llegan a
evaluarse. Los umbrales de error se mantienen intactos.

Es importante subrayar la diferencia metodológica: que un horizonte
\emph{no se pueda medir} con el rigor exigido no equivale a que el modelo
\emph{haya fallado}. Los umbrales de error (MAPE, sMAPE, errores extremos y
comparación con benchmarks) se mantienen sin relajar en todos los niveles.

\section{Caso de insumos específicos}

Las rutas que descienden hasta un insumo concreto ---por ejemplo
\textit{Concreto simple} dentro del capítulo \textit{Estructura placa huella y
drenajes}--- se procesan con exactamente el mismo flujo que las series
agregadas: no existe ningún filtro ni umbral distinto por nivel de la
jerarquía. La diferencia práctica proviene de los datos, no del método:

\begin{itemize}
  \item Un insumo con historia completa se comporta como cualquier serie agregada y admite proyección técnica en horizontes cortos.
  \item Un insumo incorporado tardíamente al anexo tiene menos observaciones válidas; al reducirse \(n\), se reduce el número de orígenes y el horizonte máximo certificable.
  \item Los insumos suelen presentar mayor volatilidad relativa que los agregados, por lo que es más frecuente que aparezcan advertencias de atípicos o de amplitud de intervalos.
\end{itemize}

Por eso la aplicación reporta siempre, junto al resultado, el número de
observaciones, el periodo cubierto, el entrenamiento inicial y el número de
ventanas de validación: son los datos que explican por qué un horizonte se
certifica, se ofrece como escenario o se niega.

\section{Alcance y limitaciones}

\begin{itemize}
  \item La calidad del resultado depende de la serie, la ruta y el periodo seleccionados.
  \item Los intervalos comunican incertidumbre estadística; no garantizan valores futuros.
  \item La aplicación no recalcula canastas, ponderaciones ni índices oficiales.
  \item Toda selección contractual y toda equivalencia requieren la justificación del profesional responsable.
\end{itemize}

\end{document}
